\section{最小缓存驻留调度}
\subsection{生命周期模型与下界}
问题一只优化节点顺序，不给缓冲区设置物理偏移，也不加入换入换出。缓冲区从申请后到释放前始终计入驻留量。对一个节点置换 $\pi$，定义处理完序列第 $k$ 个节点后的活跃指示量
\begin{equation}
y_b(k;\pi)=\mathbf 1\{\pi(a_b)\le k<\pi(f_b)\},\qquad k=0,\ldots,N.
\end{equation}
于是调度模型为
\begin{align}
\min_{\pi}\quad&H(\pi)=\max_{0\le k\le N}\sum_{b\in B}s_b y_b(k;\pi),\label{eq:peak}\\
\text{s.t.}\quad&\pi(u)<\pi(v),\qquad (u,v)\in E,\label{eq:topo}\\
&\pi(a_b)<\pi(v)<\pi(f_b),\qquad v\in V_o,\ b\in B_v,\label{eq:use}\\
&\pi:V\longrightarrow\{1,\ldots,N\}\text{ 为双射}.\nonumber
\end{align}
式\eqref{eq:use}把使用时必须活跃的语义显式列出；在输入依赖充分描述生命期时，它由原始路径蕴含。申请记正增量、释放记负增量的题面定义，与式\eqref{eq:peak}完全一致。总峰值是在同一位置对各类型求和后取最大值，不能相加各类型在不同位置达到的峰值。

每条操作的全部输入和输出缓冲区必须同时存在。因此，无论如何重排，都有
\begin{equation}
H^*\ge L_H:=\max_{v\in V_o}\sum_{b\in B_v}s_b.\label{eq:hlower}
\end{equation}
此界忽略跨操作必须保留的数据，计算成本低但可能较松。若某候选达到该界，可以证明其峰值最优；未达到时，只能据上下界判断尚未排除的改进空间。

\subsection{固定操作顺序下的生命周期收缩}
题面规定申请节点为根、释放节点为叶，六图均满足这一结构。固定非管理操作的相对顺序后，一个申请必须早于其所有显式后继以及所有引用该缓冲区的操作。称其中最早出现的节点为首次需求节点。相应地，释放必须晚于全部显式前驱和全部使用者，称其中最晚出现的节点为末次需求节点。这里不能仅以第一次读或最后一次读代替需求节点，否则可能遗漏额外管理依赖。

\textbf{收缩性质。}在上述前提下，将申请移动到首次需求节点之前，将释放移动到末次需求节点之后，不会增加驻留峰值。

证明分两步。首先，申请是根节点，向后移动不会违反其入边；又因为尚未跨过任何显式后继或使用者，其出边和使用约束仍成立。将它与中间节点逐个交换，只会删除被跨越区间内这一缓冲区的正容量贡献。其次，释放是叶节点，向前移动不会违反其出边；停止在全部前驱和使用者之后，则其入边不受影响。逐个交换同样只会减少中间区间的驻留量。其他缓冲区的申请、释放相对关系不变，因此每一步交换的最大累计量均不增加。对所有缓冲区重复此变换，得到同一操作顺序下不劣的规范化序列。

这一性质消去了大量仅在管理节点位置上不同的候选，却没有解决操作顺序的组合搜索。特别是“尽早释放”要求末次使用已结束，并不允许提前释放随后还会用到的中间数据。若推广到申请有入边或释放有出边的图，必须重新限定可移动范围，不能直接使用本性质。

\subsection{候选生成与算法对应}
在操作诱导子图中维护入度和就绪集合。每选择一个操作，先插入尚未出现且为该操作所需的申请，随后发出操作，更新剩余使用次数；只有剩余使用次数为零且全部显式前驱已完成时才插入释放。输出完毕后逐边检查式\eqref{eq:topo}，并检查节点置换及累计量首尾归零。

为了分开不同调度偏好的效果，设置三类主要操作顺序。编号顺序每次选择最小编号的就绪节点，是经过生命期收缩的强基准，而非未经处理的文件行顺序。缓存贪心先选就绪集合中编号最小的 $W$ 个节点，再最小化
\begin{equation}
q_v=n_v-\alpha d_v,\quad
n_v=\sum_{b\in B_v\setminus B_{\rm live}}s_b,\quad
d_v=\sum_{b\in B_v:\,\operatorname{remain}(b)=1}s_b.\label{eq:greedy}
\end{equation}
其中 $d_v$ 是末次使用带来的释放潜力；若还有额外释放前驱未完成，实际释放仍须等待，评分不强制兑现这一潜力。相同评分按新增量、下游路径及编号打破平局。比较 $(W,\alpha)=(32,1),(128,1),(64,2)$，窗口并非从所有就绪节点中寻找评分最好的 $W$ 项。

第三类从编号有序的终端操作出发，沿前驱作深度优先访问，以完成访问顺序得到拓扑序；前驱正序和逆序各形成一个候选。它有机会缩短单个依赖分支的持续时间，但共享缓冲区可能使分支优先与复用优先发生冲突。另保留两种流水评分产生的操作顺序参与峰值比较。最终八种顺序均经过相同生命期收缩，再按真实峰值选优，不用局部评分直接替代目标函数。

令 $Q=\sum_{v\in V_o}|B_v|$，$K=\max_v|B_v|$。图遍历及未排序的深度优先核心为 $O(N+|E|)$；实现中的编号排序另有至多 $O((N+|E|)\log N)$ 开销。窗口贪心每步扫描就绪集合，并计算至多 $W$ 个缓冲区评分，保守界为 $O(N^2\log(W+1)+NWK+|E|+Q)$。空间为 $O(N+|E|+Q)$。实际并未枚举所有拓扑序，因此计算可承受性以限制搜索范围为代价。

\subsection{六图峰值与解质量}
表\ref{tab:peaks}同时给出收缩后编号基准、选定峰值及式\eqref{eq:hlower}的下界。比值 $H_1/L_H$ 是相对一个可计算下界的比值，不是已知的最优性误差。
\begin{table}[H]\centering\caption{驻留峰值与同时使用下界（题面容量单位）}\label{tab:peaks}\small
\begin{tabular}{lrrrr}\toprule
计算图 & $H_{\rm id}$ & $H_1$ & $L_H$ & $H_1/L_H$\\\midrule
卷积小图 & 7560 & 6984 & 3144 & 2.22\\
卷积大图 & 14048 & 14048 & 932 & 15.07\\
注意力小图 & 4884 & 4884 & 512 & 9.54\\
注意力大图 & 7972 & 7972 & 512 & 15.57\\
矩阵乘小图 & 9728 & 9728 & 512 & 19.00\\
矩阵乘大图 & 35328 & 35328 & 512 & 69.00\\
\bottomrule\end{tabular}
\end{table}

卷积小图通过深度优先顺序从 7560 降至 6984，减少 576，降幅为 7.62\%。其余五图在本轮候选中未降低编号基准的峰值；选定峰值分别为 14048、4884、7972、9728 和 35328。这一结果不能归结为“五图已最优”：例如矩阵乘大图的下界仅为 512，候选与下界相距甚远，既可能包含仍可消除的跨操作生命期，也反映同时使用下界本身过弱。

局部释放评分还可能明显恶化结果。卷积小图取 $(32,1)$ 时峰值为 21762，超过编号基准的两倍。评分只看到当前操作的新增量与末次使用潜力，没有估计后续分支累计的活跃集合。这一有限视野足以解释为什么不能把局部评分的下降视为全局峰值的下降；具体多出的活跃对象尚未逐个归因，不将这一解释当作已完成的因果分解。问题二仍需独立评价物理分配与搬运量，不能只凭本问峰值筛掉其他有价值的顺序。
